\documentclass[11pt]{article}
\usepackage[a4paper,margin=2.5cm]{geometry}
\usepackage{amsmath,amssymb,amsthm}
\usepackage{hyperref}

\newtheorem{lemma}{Lemma}
\theoremstyle{definition}
\newtheorem{definition}{Definition}
\theoremstyle{remark}
\newtheorem{remark}{Remark}

\newcommand{\F}{\mathbb{F}}
\newcommand{\MS}{\mathrm{MS}}
\newcommand{\Det}{\operatorname{Det}}

\title{Multisets of $O(k)$-invariants under a labelling symmetry group\\
\large Working notes}
\author{}
\date{\today}

\begin{document}
\maketitle

\section{Setup}

Let $L = (v_1, v_2, \ldots, v_n)$ be an ordered list of $n$ vectors, each
$v_i \in \F^k$, where $\F$ is usually $\mathbb{R}$ but may be $\mathbb{C}$.

\subsection{Invariants}

\begin{definition}[Dots]
For $1 \le i \le j \le n$ let
\[
  d_{ij} = v_i \cdot v_j = v_i^{\mathsf T} v_j .
\]
There are $n + \binom{n}{2}$ of these (self-pairs included). The product is
bilinear (no complex conjugation) also when $\F = \mathbb{C}$.
\end{definition}

\begin{definition}[Epsilons]
For $1 \le i_1 < i_2 < \cdots < i_k \le n$ let
\[
  e_{i_1 i_2 \cdots i_k} = \Det\bigl[v_{i_1}, v_{i_2}, \ldots, v_{i_k}\bigr].
\]
There are $\binom{n}{k}$ of these. The inequalities are strict: any repeated
index gives a zero determinant.
\end{definition}

For non-canonically-ordered indices we use $d_{ji} = d_{ij}$ and
$e_{\ldots j \ldots i \ldots} = -e_{\ldots i \ldots j \ldots}$ (total
antisymmetry), so any index list can be brought to canonical (increasing)
order, gaining a sign $\operatorname{sgn}(\pi)$ for the sorting permutation
$\pi$ in the case of an $e$.

\subsection{The vector $w$}

Let $w$ be the column vector containing every $d$ and every $e$ in a fixed,
predetermined order. Any order is allowed, but it is convenient to keep
together the entries that form a single orbit under the permutation part of the
group $G$ (below).

\subsection{The symmetry group $G$}

$G$ is generated by
\begin{itemize}
  \item a subgroup of $S_n$ acting on the labels $1, \ldots, n$, typically a
    product $S(n_1) \times S(n_2) \times \cdots$ permuting labels only within
    blocks. For example, for $n = 5$, $S(3) \times S(2)$ permutes
    $\{v_1, v_2, v_3\}$ among themselves and $\{v_4, v_5\}$ among themselves;
  \item optionally, a parity factor $P = \{1, p\}$, so that
    $|G \times P| = 2|G|$.
\end{itemize}

\begin{definition}[Parity]
$p$ acts on every vector by the same fixed reflection,
\[
  p : v_i \mapsto R\, v_i, \qquad R = \operatorname{diag}(1, \ldots, 1, -1),
\]
that is, it flips the sign of the last component of every vector.
\end{definition}

\begin{lemma}\label{lem:parity}
For every $k$, $p$ leaves each $d$ unchanged and negates each $e$:
\[
  p : d_{ij} \mapsto d_{ij}, \qquad p : e_{i_1 \cdots i_k} \mapsto -e_{i_1 \cdots i_k}.
\]
\end{lemma}
\begin{proof}
$R^{\mathsf T} R = I$, so $(R v_i)^{\mathsf T}(R v_j) = v_i^{\mathsf T} v_j$.
Also $\Det[R v_{i_1}, \ldots, R v_{i_k}] = \det(R)\,\Det[v_{i_1}, \ldots, v_{i_k}]$
and $\det R = -1$.
\end{proof}

\begin{remark}
The proof uses only $R^{\mathsf T} R = I$ and $\det R = -1$. So any reflection,
or more generally any orthogonal matrix with $\det = -1$, acts identically on
$w$, and the choice of which component to flip is immaterial. Because $R$ is
real, the $d$'s are invariant under either a bilinear or a Hermitian product.
The $e$'s, however, pair naturally with the bilinear $d$.
\end{remark}

\begin{remark}[Point inversion is \emph{not} the right parity]
The alternative $v \mapsto -v$ (that is, $R = -I_k$) has
$\det(-I_k) = (-1)^k$. It negates the $e$'s only for odd $k$. For even $k$
(for example $k = 2$, where $-I_2$ is a rotation by $\pi$) it leaves $w$
completely unchanged. This is why parity is defined as a reflection above.
\end{remark}

\begin{remark}[Indefinite metrics, for example Minkowski space]\label{rem:metric}
The dot products may be taken with any non-degenerate diagonal metric,
$d_{ij} = v_i^{\mathsf T} \eta\, v_j$ with $\eta = \operatorname{diag}(\pm 1,
\ldots, \pm 1)$ (symembed's option \texttt{metric}). The determinants $e$ do not
involve the metric. Nothing downstream of $w$ depends on it either: every
construction of Sections~3 and~4, every output count, and every statement
about injectivity on $\MS$ is unchanged. Parity is then any isometry of
determinant $-1$, and Lemma~1 holds as before.

What does change is the last step, from $\MS$ back to $L$. Over a definite
metric the Gram matrix always determines $L$ up to $O(k)$. Over an indefinite
one that is true only when the $v_i$ span a \emph{non-degenerate} subspace.
Then equal Gram matrices force equal linear relations among the vectors, so the
correspondence between the two configurations is an isometry between their
spans, and Witt's theorem extends it to the whole space. (The $e$'s fix the
determinant, as before.) The step fails for degenerate spans, which need null
vectors. The simplest case is a null vector $n$ and the zero vector, which have
the same (zero) Gram matrix.

\emph{Example.} In $1+1$ dimensions, with $P$ in $G$ so that only dot products
are needed, take
\[
  L_1 = \bigl((3,3), (2,2)\bigr), \qquad L_2 = \bigl((3,3), (1,1)\bigr), \qquad
  L_3 = \bigl((3,3), (0,0)\bigr), \qquad L_0 = \bigl((0,0), (0,0)\bigr).
\]
Every dot product vanishes in all four. In light-cone coordinates $u = t + x$,
$v = t - x$, the non-zero vectors lie on the ray $v = 0$, with
$u = (6,4)$, $(6,2)$ and $(6,0)$. A boost multiplies every $u$ by the same
factor $e^\eta$, so the ratio $u_1 : u_2$ is invariant. Parity exchanges $u$
and $v$, time reversal makes vectors past-pointing, and relabelling only swaps
the ratio. So $L_1$, $L_2$, $L_3$ and $L_0$ lie in four different orbits, even
up to~$G$. Yet boosting with $\eta \to -\infty$ shrinks every vector towards
zero, so all four orbit closures contain $L_0$, and every continuous invariant
takes the same value on all four. The quantity that does tell $L_1$ from $L_2$,
the momentum fraction $z = u_1/(u_1+u_2)$ ($0.6$ against $0.75$), is
Lorentz-invariant on exactly-collinear massless pairs but cannot be extended
continuously. Near $L_0$, which both orbits approach, it takes every value.
This rhymes with collinear safety in physics: collinear-safe observables are
precisely those that do not depend on such splitting fractions.
\begin{itemize}
  \item \emph{No continuous invariant can do better.} The Lorentz group is not
    compact, and its orbits need not be closed. Boosts carry $n$ arbitrarily
    close to~$0$, so any continuous invariant takes the same value on $n$ as
    on~$0$. Every continuous invariant is constant on orbit
    \emph{closures}, so the most any continuous embedding can do is separate
    orbit closures.
  \item \emph{That is also all that is needed.} The purpose of the embedding
    is that every continuous $G$-invariant function of $L$ should be a function
    of the outputs. Continuous invariant functions are themselves constant on
    orbit closures. So failing to tell apart configurations whose orbit
    closures meet costs nothing.
  \item \emph{Physical momenta.} For future-pointing timelike or null
    vectors, a span containing a timelike vector is non-degenerate, and so is
    the span of two non-parallel null vectors. So the only degenerate
    configurations are those in which every vector is null and all are
    mutually parallel (or zero). Their orbit closures contain~$0$. Elsewhere,
    the $d$'s and $e$'s determine $L$ up to $SO(1,3)$, and, because the vectors
    are future-pointing, up to the proper orthochronous group. Separately,
    note that the $d$'s and $e$'s cannot detect time reversal: if past-pointing
    vectors were allowed, invariants of the orthochronous group such as
    $\operatorname{sign}(p^0) \max(p^2, 0)$ would not be functions of them.
\end{itemize}
For a definite metric every orbit is closed, so none of this arises.\footnote{%
  \emph{Note to self, on \cite{GHL20} and \cite{LHG22}.} These papers use the
  fact that the Lorentz- and parity-invariant polynomials are generated by the
  dot products (Weyl), together with the density of polynomials, to conclude
  that every invariant function is a function of the dots. That step needs care
  for the non-compact Lorentz group. Density of polynomials does not by itself
  give density of \emph{invariant} polynomials among invariant functions: there
  is no averaging over a non-compact group. And the conclusion fails if
  past-pointing momenta are allowed (the $\operatorname{sign}(p^0)$ example
  above). Moreover, the dots do not determine the momenta up to Lorentz
  transformations for collinear massless configurations. The conclusions do
  survive for \emph{continuous} functions of \emph{future-pointing} momenta:
  continuous invariants are constant on orbit closures, and on that domain the
  dots separate orbit closures. That justification, rather than the
  density argument alone, is the one to give.}
\end{remark}

\subsection{Group action on $w$ and the multiset}

For $g \in G$, the vector $g \cdot w$ is obtained as follows:
\begin{enumerate}
  \item relabel the indices of each entry according to the permutation part
    of $g$;
  \item restore canonical index order, using $d_{ji} = d_{ij}$ and the
    antisymmetry sign for each $e$;
  \item if $g$ contains $p$, negate every $e$ (Lemma~\ref{lem:parity}).
\end{enumerate}
The resulting entries are re-expressed in terms of the original (canonical)
$d$'s and $e$'s, keeping the slot positions of $w$ fixed. Define the multiset
\[
  \MS = \{\, g \cdot w \mid g \in G \,\},
\]
with multiplicities, so that $|\MS| = |G|$ counted with multiplicity.

\section{Worked example: $k = 2$, $L = (a, b, p)$}

Here $n = 3$. To avoid a clash with the parity element, the vector named $p$ in
$L$ is written as a subscript only. There are $3 + 3 = 6$ dots and
$\binom{3}{2} = 3$ epsilons. Ordered by orbits of $S(2) \times S(1)$ (bars
separate the orbits):
\[
  w = \bigl(\, d_{aa}, d_{bb} \;\big|\; d_{pp} \;\big|\; d_{ab} \;\big|\;
      d_{ap}, d_{bp} \;\big|\; e_{ab} \;\big|\; e_{ap}, e_{bp} \,\bigr)^{\mathsf T}.
\]
Let $\sigma = (ab)$. Under $\sigma$: $d_{ab} \mapsto d_{ba} = d_{ab}$,
$e_{ab} \mapsto e_{ba} = -e_{ab}$, and $e_{ap} \leftrightarrow e_{bp}$ with no
sign change.

\subsection{$G = S(2) \times S(1)$}

\[
  \MS = \left\{\;
  \underbrace{\begin{pmatrix}
    d_{aa}\\ d_{bb}\\ d_{pp}\\ d_{ab}\\ d_{ap}\\ d_{bp}\\ e_{ab}\\ e_{ap}\\ e_{bp}
  \end{pmatrix}}_{1},\;
  \underbrace{\begin{pmatrix}
    d_{bb}\\ d_{aa}\\ d_{pp}\\ d_{ab}\\ d_{bp}\\ d_{ap}\\ -e_{ab}\\ e_{bp}\\ e_{ap}
  \end{pmatrix}}_{\sigma}
  \;\right\}
\]

\subsection{$G = S(2) \times S(1) \times P$}

\[
  \MS = \left\{\;
  \underbrace{\begin{pmatrix}
    d_{aa}\\ d_{bb}\\ d_{pp}\\ d_{ab}\\ d_{ap}\\ d_{bp}\\ e_{ab}\\ e_{ap}\\ e_{bp}
  \end{pmatrix}}_{1},\;
  \underbrace{\begin{pmatrix}
    d_{bb}\\ d_{aa}\\ d_{pp}\\ d_{ab}\\ d_{bp}\\ d_{ap}\\ -e_{ab}\\ e_{bp}\\ e_{ap}
  \end{pmatrix}}_{\sigma},\;
  \underbrace{\begin{pmatrix}
    d_{aa}\\ d_{bb}\\ d_{pp}\\ d_{ab}\\ d_{ap}\\ d_{bp}\\ -e_{ab}\\ -e_{ap}\\ -e_{bp}
  \end{pmatrix}}_{p},\;
  \underbrace{\begin{pmatrix}
    d_{bb}\\ d_{aa}\\ d_{pp}\\ d_{ab}\\ d_{bp}\\ d_{ap}\\ e_{ab}\\ -e_{bp}\\ -e_{ap}
  \end{pmatrix}}_{\sigma p}
  \;\right\}
\]
In $\sigma p$, the swap sign and the parity sign cancel on $e_{ab}$.

\section{Layer 2: embedding the multiset}

The multiset $\MS$ is not yet a usable output. Next we apply a
continuous, injective and algebraic \emph{multiset embedding}
\[
  \Phi : \bigl(\F^m\bigr)^{|G|} / S_{|G|} \longrightarrow \F^D, \qquad m = \dim w,
\]
and express each component of $\Phi(\MS)$ explicitly as a polynomial in
the $d$'s and $e$'s. Several choices of $\Phi$ are possible. We start with the
Bursarial embedding because its outputs are explicit low-degree polynomials,
which makes it the easiest to work with.

\subsection{The Bursarial embedding}

Index the slots of $w$ from $0$ to $m-1$, so $w = (w_0, \ldots, w_{m-1})^{\mathsf T}$.
For each $g \in G$ define
\[
  q_g(x) = \sum_{j=0}^{m-1} (g \cdot w)_j \, x^j ,
\]
and multiply one linear factor in $y$ for each group element:
\[
  B(x, y) = \prod_{g \in G} \bigl( y + q_g(x) \bigr).
\]

\paragraph{Elementary symmetric polynomials.}
For $N$ quantities $z_1, \ldots, z_N$, the \emph{$c$-th elementary symmetric
polynomial} is the sum of all products of $c$ of them with distinct indices:
\[
  E_c(z_1, \ldots, z_N) = \sum_{1 \le i_1 < i_2 < \cdots < i_c \le N}
                          z_{i_1} z_{i_2} \cdots z_{i_c},
  \qquad E_0 = 1.
\]
For example, $E_1 = z_1 + \cdots + z_N$ and $E_N = z_1 z_2 \cdots z_N$. When
$N = 3$, $E_2 = z_1 z_2 + z_1 z_3 + z_2 z_3$. They are the coefficients that
appear when a product of linear factors is multiplied out:
\[
  \prod_{i=1}^{N} (y + z_i) = \sum_{c=0}^{N} y^{N-c}\, E_c(z_1, \ldots, z_N).
\]
(For $N = 2$ this is $(y + z_1)(y + z_2) = y^2 + y(z_1 + z_2) + z_1 z_2$.)

\paragraph{Applying this to $B$.}
Take $N = |G|$ and let the $z$'s be the polynomials $q_g(x)$, one for each
$g \in G$. Write
\[
  E_c(x) \equiv E_c\bigl(\{q_g(x)\}_{g \in G}\bigr)
        = \sum_{\substack{\{g_1, \ldots, g_c\} \subseteq G\\ \text{($c$ distinct elements)}}}
          q_{g_1}(x)\, q_{g_2}(x) \cdots q_{g_c}(x).
\]
Then
\[
  B(x, y) = \sum_{c=0}^{|G|} y^{|G|-c}\, E_c(x).
\]
Each $q_g$ has degree (at most) $m-1$ in $x$, so $E_c(x)$, a sum of products of
$c$ of them, is a polynomial in $x$ of degree (at most) $c(m-1)$. Write its
coefficients as $E_{c,j}$:
\[
  E_c(x) = \sum_{j=0}^{c(m-1)} E_{c,j}\, x^j .
\]
(The notation $[x^j]\,f$ is sometimes used for the coefficient of $x^j$ in
$f$, so that $E_{c,j} = [x^j]\,E_c$.) The leading term $y^{|G|}$ ($c = 0$)
carries no information. The embedding is the list of all the remaining
coefficients:
\[
  \Phi(\MS) = \bigl( E_{c,j} \bigr)_{1 \le c \le |G|,\; 0 \le j \le c(m-1)},
\]
which has
\[
  D = \sum_{c=1}^{|G|} \bigl( c(m-1) + 1 \bigr) = |G| + (m-1)\,\frac{|G|(|G|+1)}{2}
\]
components. This is the construction used in Echinocoder's
\path{Cinf_sympy_bursar_embedder_for_array_of_reals_as_multiset.py},
with $n \to |G|$ and $k \to m$.

\paragraph{Properties.}
\begin{itemize}
  \item \emph{Injective on multisets.} $\F[x][y]$ is a unique factorisation
    domain, and each factor $y + q_g(x)$ is monic in $y$. So $B$ determines
    the multiset $\{q_g\}$, and hence $\MS$.
  \item \emph{Invariant.} $B$ is a product over all of $G$, so $\Phi(\MS)$ is
    unchanged when $L$ is acted on by any $g \in G$.
  \item \emph{Polynomial.} $E_{c,j}$ is a homogeneous polynomial of
    degree $c$ in the entries of $w$, that is, in the $d$'s and $e$'s:
    \[
      E_{c,j} = \sum_{\substack{\{g_1, \ldots, g_c\} \subseteq G}} \;
      \sum_{\substack{j_1 + \cdots + j_c = j}} \;
      \prod_{r=1}^{c} (g_r \cdot w)_{j_r}.
    \]
  \item \emph{Order-dependent coordinates.} The slot order of $w$ decides
    which entries multiply which powers of $x$. Changing it changes the
    coordinates of $\Phi$, but not its injectivity.
\end{itemize}

\paragraph{What injectivity is actually required.}
We never want to recover $L$ itself. The purpose of $\MS$ is to forget exactly
the information that $G$ moves. Using $d$'s and $e$'s erases the orientation
of $L$ in $\F^k$, that is, the action of $SO(k)$. Over $\mathbb{R}$ the Gram
matrix $(d_{ij})$ determines $L$ up to $O(k)$. The $e$'s then fix the handedness
whenever the $v_i$ span $\F^k$; when they do not, there is no handedness to fix.
Each factor of $G$ erases something further. For example, $S(3)$ erases which of
$a, b, c$ is which, and $P$, if present, erases handedness, so that $L$ is then
known only up to $O(k)$. In that case the $e$'s carry nothing beyond the $d$'s.
Every $G$-invariant is then even in the $e$'s, and by Cauchy--Binet
$e_I\, e_J = \det\bigl[d_{i_r j_s}\bigr]_{r,s=1}^{k}$. They are kept in $w$
here only for uniformity. A practical implementation would leave them out
whenever $P \subseteq G$, because they greatly enlarge $\MS$ and everything that
then has to be pruned.
(Over $\mathbb{C}$ with a bilinear product the degenerate cases, such as
nonzero $v$ with $v \cdot v = 0$, still need checking.) Two notions of
injectivity should therefore be kept apart:
\begin{enumerate}
  \item[(A)] \emph{Formal injectivity.} Treat the $m$ entries of $w$ as
    independent variables, and require $\Phi$ to separate $\{g \cdot w\}$ for
    \emph{every} $w \in \F^m$. This is all the Bursarial embedding knows about.
  \item[(B)] \emph{Injectivity on the multisets that actually occur.} Require
    \[
      \Phi\bigl(\MS(L)\bigr) = \Phi\bigl(\MS(L')\bigr)
      \;\Longrightarrow\; \MS(L) = \MS(L'),
    \]
    equivalently $L' = g \cdot (h\, L)$ for some $g \in G$ and orthogonal $h$
    (special orthogonal if $P \not\subseteq G$). Only those $w = w(L)$ lying
    on the variety cut out by the Gram and Plücker relations need to be
    separated.
\end{enumerate}
(B) is the target. (A) implies (B), but (B) asks less of $\Phi$, so it should
be achievable with fewer outputs.

\paragraph{Redundancy.}
$\Phi$ is deliberately oversized, for two separate reasons. First, it does not
know that $\MS$ is structured. $\MS$ is the $G$-orbit of a single point $w$,
and each slot of $g \cdot w$ is $\pm$ an entry from the same $G$-orbit block of
$w$, so many coefficients vanish identically and many others are functionally
dependent on the rest. Second, it does not know the Gram and Plücker
relations among the $d$'s and $e$'s, which are what separate target (B) from
target (A). Exploiting both kinds of redundancy is the next layer. The present goal is
only the oversized expressions, which are certainly injective, generated
algorithmically and verifiably for general $G$.

\subsection{Example: $k = 2$, $L = (a, b, p)$, $G = S(2) \times S(1)$}

With the slot order of Section~2, $m = 9$, $|G| = 2$ and $D = 2 + 8 \cdot 3 = 26$:
\begin{align*}
  q_1(x)      &= d_{aa} + d_{bb} x + d_{pp} x^2 + d_{ab} x^3 + d_{ap} x^4 + d_{bp} x^5
                 + e_{ab} x^6 + e_{ap} x^7 + e_{bp} x^8, \\
  q_\sigma(x) &= d_{bb} + d_{aa} x + d_{pp} x^2 + d_{ab} x^3 + d_{bp} x^4 + d_{ap} x^5
                 - e_{ab} x^6 + e_{bp} x^7 + e_{ap} x^8 .
\end{align*}
The $c = 1$ block ($E_1 = q_1 + q_\sigma$, 9 coefficients $E_{1,0}, \ldots, E_{1,8}$) is
\[
  \bigl(\, d_{aa}{+}d_{bb},\; d_{aa}{+}d_{bb},\; 2d_{pp},\; 2d_{ab},\;
     d_{ap}{+}d_{bp},\; d_{ap}{+}d_{bp},\; 0,\; e_{ap}{+}e_{bp},\; e_{ap}{+}e_{bp} \,\bigr),
\]
which already shows a vanishing coefficient and repeated ones. The $c = 2$
block ($E_2 = q_1 q_\sigma$, 17 coefficients $E_{2,0}, \ldots, E_{2,16}$) contains, for example,
\[
  E_{2,12} = 2\bigl( d_{ap} e_{ap} + d_{bp} e_{bp} \bigr) - e_{ab}^2 .
\]
This is where the information about $e_{ab}$ lost from $E_1$ reappears.

\subsection{Involution packing}\label{sec:involution}

The Bursarial embedding can be made roughly half the size, while staying
injective by construction (Level~0), whenever the multiset is invariant under a
fixed linear involution. For $\MS$ this is always the case if $G$ contains an
element acting on $w$ as a non-trivial involution.

\paragraph{The general mechanism.}
Let $T$ be a real $m \times m$ matrix with $T^2 = I$, and let $M$ be a multiset
of vectors in $\mathbb{R}^m$ that is invariant under $T$ as a multiset, that
is, $\{T v : v \in M\} = M$. Choose integer bases
$f_1, \ldots, f_{m_+}$ and $h_1, \ldots, h_{m_-}$ of linear functionals with
$f_a \circ T = f_a$ and $h_b \circ T = -h_b$. Here $m_\pm$ are the dimensions of
the $\pm 1$ eigenspaces, $m_+ + m_- = m$, and $m_+ - m_- = \operatorname{tr} T$.
Pad the shorter list with zero functionals to a common length
$r = \max(m_+, m_-) = (m + |\operatorname{tr} T|)/2$, and pack each vector as
\[
  \phi(v) = \bigl( f_1(v) + i\,h_1(v),\ \ldots,\ f_r(v) + i\,h_r(v) \bigr) \in \mathbb{C}^r .
\]
Then $\phi$ is injective, and $\phi(Tv) = \overline{\phi(v)}$, entry by entry. With
$q_v(x) = \sum_j \phi(v)_j\, x^j$, the complex Bursarial polynomial
\[
  B(x, y) = \prod_{v \in M} \bigl( y + q_v(x) \bigr)
\]
therefore satisfies
\[
  \overline{B} = \prod_{v \in M} \bigl( y + q_{Tv}(x) \bigr) = B ,
\]
because $v \mapsto Tv$ permutes $M$. So every coefficient of $B$ is real.
$\mathbb{C}[x][y]$ is a unique factorisation domain and each factor is monic in
$y$, so $B$ determines the multiset $\{q_v\}$, hence $\{\phi(v)\}$, hence $M$.
The real coefficients of $B$ are therefore an injective embedding of $M$ with
\[
  D = |M| + (r - 1)\,\frac{|M|(|M| + 1)}{2}
\]
outputs, instead of $|M| + (m - 1)|M|(|M| + 1)/2$.

\paragraph{Application to $\MS$.}
Let $g_0 \in G$ act on $w$ as a non-trivial involution, and take $T$ to be that
action. Because $g \mapsto g_0 g$ permutes $G$,
\[
  T(g \cdot w) = (g_0 g) \cdot w ,
\]
so $\MS$ is invariant under $T$ and the mechanism applies. The members of
$\MS$ pair up as $(g \cdot w,\ g_0 g \cdot w)$. Note that $g_0$ multiplies on the
\emph{left}. Pairing $g \cdot w$ with $g g_0 \cdot w$ instead would relate the
two by $g g_0 g^{-1}$, which depends on $g$ and is not a single fixed $T$. When
$T$ is a signed permutation, $(Tw)_s = \sigma_s\, w_{\pi(s)}$, integer
functionals can be read off slot by slot:
\begin{itemize}
  \item a fixed slot ($\pi(s) = s$) gives $w_s$, in the $+1$ part if
    $\sigma_s = +1$ and in the $-1$ part otherwise;
  \item a swapped pair ($\pi(s) = t \ne s$) gives $w_s + \sigma_s w_t$ in the
    $+1$ part and $w_s - \sigma_s w_t$ in the $-1$ part.
\end{itemize}
The best $g_0$ is the one with smallest $|\operatorname{tr} T|$. The special case
$g_0 = p$ (parity) packs every $d$ with an $e$. The case of $g_0$ swapping the
two members of a zipped pair is symcoder's $z \mapsto i\bar z$ compression.

\paragraph{Example.}
For $S(2)$ acting on $(a, b)$ with $k = 2$, take $g_0 = (ab)$, which has trace~0.
Then
\[
  z_1 = (d_{aa} + d_{bb}) + i\,(d_{aa} - d_{bb}), \qquad
  z_2 = d_{ab} + i\,e_{ab} = \bar z_a z_b ,
\]
and the 5 outputs are
\[
  2(d_{aa} + d_{bb}),\quad 2 d_{ab},\quad 2(d_{aa}^2 + d_{bb}^2),\quad
  2\bigl(d_{ab}(d_{aa} + d_{bb}) + e_{ab}(d_{aa} - d_{bb})\bigr),\quad
  d_{ab}^2 + e_{ab}^2 .
\]
The unpacked embedding has 11.

\paragraph{Sizes.}
The table gives output counts before any reduction.
\begin{center}
\begin{tabular}{lrrrr}
  case & $g_0$ & $\operatorname{tr}$ & plain & packed \\ \hline
  $S(2)$, $k=2$                 & $(ab)$        & $0$  & 11     & 5 \\
  $S(3)$, $k=2$                 & $(bc)$        & $+1$ & 174    & 90 \\
  $S(2)\times S(1)$, $k=2$      & $(ab)$        & $+1$ & 26     & 14 \\
  $S(5)$, $k=4$                 & $(de)$        & $+4$ & 138060 & 79980 \\
  $S(3)\times S(2)$, $k=3$      & $(bc)(pq)$    & $+1$ & 1884   & 948 \\
  $S(5)\times P$, $k=4$         & $(bc)(de)p$   & $+2$ & 549720 & 289440
\end{tabular}
\end{center}

\paragraph{Limits.}
Complex conjugation is the only non-trivial symmetry of $\mathbb{C}$ over
$\mathbb{R}$, so conjugation alone saves at most a factor of about~2. (Using
two commuting involutions in the same way would need an algebra with two
independent conjugations. Bicomplex numbers have zero divisors, which breaks
unique factorisation, and quaternions do not commute, so the product depends on
the order of its factors.) This is \emph{not} a limit on exploiting the group,
however. Section~\ref{sec:graded} gets much larger factors by also letting a
group element act on the variable~$x$.

On the simple cases of Section~\ref{sec:exploration}, full polynomial closure
(in $L$) leaves exactly as many outputs from the packed embedding as from the
plain one. The packing does not change what survives that (expensive)
reduction. What it gives is the saving for free, with no reduction work, and it
does better than the plain embedding followed by the linear (safe) reductions.

\subsection{Graded packing}\label{sec:graded}

Involution packing let one group element act on the \emph{coefficients} of
$B$, as complex conjugation. Graded packing also lets a group element act on the
\emph{variable} $x$. That kills most coefficients outright, and the saving grows
with the order of the element.

\paragraph{A warm-up: even polynomials.}
Suppose a linear map $T$ permutes the members of a multiset $M$, and acts on
each vector by flipping the sign of its odd-numbered entries:
$(Tv)_j = (-1)^j v_j$. Then
\[
  q_{Tv}(x) = \sum_j (-1)^j v_j x^j = q_v(-x),
\]
and so
\[
  B(-x, y) = \prod_{v \in M} \bigl( y + q_v(-x) \bigr)
           = \prod_{v \in M} \bigl( y + q_{Tv}(x) \bigr) = B(x, y),
\]
because $v \mapsto Tv$ only reorders the factors. So $B$ is an even function of
$x$, and every coefficient of an odd power of $x$ is zero for every input.
Those coefficients never need to be stored. This is the same factor of~2 as
involution packing, obtained with real numbers by putting the $+1$ coordinates
of $T$ at even positions and the $-1$ coordinates at odd positions.

\paragraph{Roots of unity.}
Now replace $-1$ by $\zeta = e^{2\pi i/r}$. If a symmetry $h$ of $M$
multiplies the coordinate at position $j$ by $\zeta^j$, then
$q_{hv}(x) = q_v(\zeta x)$, so $B(\zeta x, y) = B(x, y)$. Comparing the
coefficients of $y^{|M|-c} x^j$ on both sides gives
$E_{c,j} = \zeta^j E_{c,j}$, and $\zeta^j \ne 1$ unless $r$ divides $j$.
Hence
\[
  E_{c,j} = 0 \quad \text{unless } r \mid j ,
\]
so only about one coefficient in $r$ survives.

\paragraph{Finding such coordinates.}
For $\MS$, take any $h \in G$, and let $r$ be the order of its action on $w$.
Since $h^r = 1$, every eigenvalue of that action is an $r$-th root of unity
$\zeta^k$. Call $k \in \{0, \ldots, r-1\}$ the \emph{weight}. Use
eigen-coordinates: linear functionals $p$ with $p(h \cdot v) = \zeta^k p(v)$.
They are real when $\zeta^k = \pm 1$, and complex otherwise. Put each coordinate
of weight $k$ at a position $j \equiv k \pmod r$, and leave unused positions
empty (zero). A complex eigen-coordinate of weight $k$ and its conjugate
(weight $-k$) carry the same information, so only one of each pair is kept. It
can go at residue $k$ or at residue $r-k$, whichever keeps the residues
balanced. Balance matters, because the length of $q$ is set by the busiest
residue.

\paragraph{Making the coefficients real.}
The surviving coefficients are complex in general, and would cost two reals
each. But suppose $G$ also contains an involution $s$ with $s h s = h^{-1}$.
Then $h$ and $s$ generate a dihedral group, like the rotations ($h$) and
reflections ($s$) of a regular $r$-gon. The coordinates can then be chosen so
that also
\[
  p(s \cdot v) = \overline{p(v)} .
\]
(Both $s$ and complex conjugation map the weight-$k$ space to the weight-$(-k)$
space. Their composite is a ``real structure'' on each weight space, and
coordinates fixed by it have exactly this property.) For the real weight
spaces ($k = 0$, and $k = r/2$ if $r$ is even) this just means pairing an
$s$-even coordinate with an $s$-odd one as (even)${}+i\,$(odd), as in involution
packing. Then $\overline{B} = B$, as in Section~\ref{sec:involution}, and every
surviving coefficient is real.

\paragraph{Invariants, and the scalar mode.}
Coordinates fixed by every element of $G$, such as $d_{aa} + d_{bb}$, are equal
for every member of $\MS$. They can be output directly and left out of the
packing. Occasionally $h$ then multiplies \emph{every} remaining coordinate by
the same factor $\zeta^\rho$. In that case the coordinates are placed
consecutively and the grading moves to $y$: from $q_{hv} = \zeta^\rho q_v$ it
follows that $E_c = \zeta^{\rho c} E_c$, so whole blocks $E_c$ vanish unless $r$
divides $\rho c$.

\paragraph{Exact zeros from symmetries of single coordinates.}
The grading is one instance of a more general argument, which also finds
coefficients that vanish for subtler reasons. Expanding the product,
\[
  E_{c,j} = \sum_{J} F_J ,
\]
where $J$ runs over \emph{patterns}: multisets $\{a_1, \ldots, a_c\}$ of
occupied positions with $a_1 + \cdots + a_c = j$. Here $F_J$ collects the
products $\prod_{g \in S} \phi_{a_g}(g \cdot w)$ over all $c$-element subsets
$S \subseteq G$ and assignments $g \mapsto a_g$ whose positions form the
multiset~$J$. So $F_J$ depends only on the multiset of tuples
$\{(\phi_a(g \cdot w))_{a \in J}\}_{g \in G}$, and it has degree $n_a$ in the
$a$-th entries, where $n_a$ is the number of times $a$ occurs in $J$.

Now suppose some $t \in G$ multiplies each coordinate in the pattern by a
scalar: $\phi_a(t \cdot v) = \chi_a\, \phi_a(v)$ for every $a \in J$. Because
$g \mapsto tg$ permutes $G$, that multiset of tuples is unchanged when every
$a$-th entry is multiplied by $\chi_a$. So
\[
  F_J = \Bigl(\prod_{a} \chi_a^{\,n_a}\Bigr) F_J ,
\]
and $F_J = 0$ unless $\prod_a \chi_a^{n_a} = 1$. If every pattern of $E_{c,j}$
is killed in this way, then $E_{c,j}$ vanishes for every input.
\begin{itemize}
  \item Taking $t = h$, which scales every coordinate, gives back the grading.
  \item A coefficient with no patterns at all is an unreachable power.
\end{itemize}
The code finds, exactly in $\mathbb{Z}[\zeta_L]$, which group elements scale
which coordinates. It then examines every coefficient within a small distance of
either end of $E_c(x)$, where the patterns are few because most positions must
equal the end position.

\emph{Example:} for $S(3)\times S(2)$ with $k = 3$, position~0 holds the purely
imaginary coordinate $i\,e_{abc}$. The reflection $s$ multiplies it by $-1$, so
every pattern of $E_{c,0}$, which can only be $\{0, \ldots, 0\}$, picks up the
factor $(-1)^c$. Hence $E_{c,0} = 0$ for every odd $c$. Equivalently, over the
12 group elements the coordinate takes the values $\pm i\,e_{abc}$ six times
each, and $E_{c,0}$ is a coefficient of $(y^2 + e_{abc}^2)^6$. Similar
cancellations occur near the top power. In all, 12 coefficients vanish for this
reason, taking that case from 203 to 191 outputs. These include four that the
Level~1 searches had found individually. ($S(3)\times S(2)$ with $k = 4$ goes
from 130 to 124.)

\paragraph{Why this is Level 0.}
The packing (together with any invariants output directly) is an injective
linear map from $w$. The complex Bursarial embedding is injective on multisets
of complex vectors, by unique factorisation in $\mathbb{C}[x][y]$. The omitted
coefficients are zero for every input, so they carry no information, and when
$s$ is used the imaginary parts of the kept coefficients are zero too. (Besides
the powers killed by the grading, the code also omits any $E_{c,j}$ whose $j$
cannot be written as a sum of $c$ \emph{occupied} positions. When some positions
are empty padding, these are zero for the simplest of reasons: the product of
$c$ polynomials $q_g$ contains no other powers of~$x$.)
Everything is exact. In practice the code tries one $h$ from each conjugacy
class, every suitable $s$ (or none), and with or without taking the invariants
out, and keeps whichever choice gives the fewest outputs. The choice $h = e$
recovers involution packing.

\paragraph{Example: $S(2)$, $(a, b)$, $k = 2$.}
Here $w = (d_{aa}, d_{bb}, d_{ab}, e_{ab})$. The invariants $d_{aa} + d_{bb}$ and
$d_{ab}$ are output directly. The swap $h = (ab)$ multiplies both remaining
coordinates, $d_{aa} - d_{bb}$ and $e_{ab}$, by $-1$, so this is the scalar mode
with $r = 2$. Packing them as $z = (d_{aa} - d_{bb}) + i\, e_{ab}$, the multiset is
$\{z, -z\}$ and $B = (y + z)(y - z) = y^2 - z^2$. So $E_1 = 0$, and the only
coefficient is $-z^2$. It is complex, because no element of $G$ acts on $z$ as
conjugation. The four outputs are
\[
  d_{aa} + d_{bb}, \qquad d_{ab}, \qquad
  \mathrm{Re}\, z^2 = (d_{aa} - d_{bb})^2 - e_{ab}^2, \qquad
  \mathrm{Im}\, z^2 = 2\, e_{ab} (d_{aa} - d_{bb}) .
\]
This is the hand-made zip of Section~\ref{sec:exploration}, and it equals the
count left by the full polynomial closure. The plain embedding has 11 outputs.

\paragraph{Example: $S(3)$, $(a, b, c)$, $k = 2$.}
Take $h = (abc)$ ($r = 3$) and $s = (bc)$. The 9 slots of $w$ give three
weight-0 coordinates and three complex coordinates of weight $\pm 1$.
\begin{itemize}
  \item The weight-0 coordinates are $\sum d_{ii}$ and $\sum_{i<j} d_{ij}$,
    which are $s$-even, and one combination of the $e$'s, which is $s$-odd.
    Paired, they become 2 complex coordinates, placed at positions 0 and 3.
  \item The three complex coordinates are split between residues 1 and 2, at
    positions 1, 4 and 2.
\end{itemize}
So each $q$ has length $W = 5$, and $E_c$ has degree $4c$. Keeping only powers
divisible by 3 leaves $\lfloor 4c/3 \rfloor + 1$ real coefficients for each
$c = 1, \ldots, 6$, that is $2 + 3 + 5 + 6 + 7 + 9 = 32$, against 174 for the
plain embedding.

\paragraph{Sizes.}
In the table below, the ``plain'' and ``involution'' columns count every
coefficient of those constructions. The ``graded'' column counts the outputs of
the graded construction itself. Coefficients that vanish identically (by the
grading, as unreachable powers, or by character symmetries) are part of that
construction, and are already omitted. The last column also applies the
optional exact Level~1 reductions (\texttt{reductions="safe"}).
\begin{center}
\small
\begin{tabular}{llrrrrr}
  case & $(h, s)$ & plain & involution & graded & factor & ${}+{}$safe \\ \hline
  $S(2)$, $k=2$            & $(ab)$, scalar          & 11     & 5     & 4     & 2.8  & 4 \\
  $S(2)\times S(1)$, $k=2$ & $(ab)$, scalar          & 26     & 14    & 11    & 2.4  & 11 \\
  $S(3)$, $k=2$            & $(abc)$, $(bc)$         & 174    & 90    & 32    & 5.4  & 32 \\
  $S(4)$, $k=2$            & $(bcd)$, $(cd)$         & 4524   &       & 718   & 6.3  & 715 \\
  $S(4)$, $k=3$            & $(abcd)$, $(bd)$        & 3924   &       & 469   & 8.4  & 468 \\
  $S(3)\times S(2)$, $k=3$ & $(abc)(pq)$, $(bc)$     & 1884   & 948   & 191   & 9.9  & 191 \\
  $S(3)\times S(2)$, $k=4$ & $(abc)(pq)$, $(bc)$     & 1494   &       & 124   & 12.0 & 124 \\
  $S(5)$, $k=4$            & $(abcde)$, $(be)(cd)$   & 138060 & 79980 & 13141 & 10.5 & 13139
\end{tabular}
\end{center}

\paragraph{Limits.}
\begin{itemize}
  \item \emph{The factor is bounded by the dihedral subgroup used.} The saving
    is roughly $|\langle h, s\rangle| = 2r$, and $r$ is at most the largest
    order of an element of $G$ (6 for $S(5)$, 30 for $S(10)$). That is tiny
    compared with $|G| = n!$. This is a large constant factor for moderate
    $n$, not a change in the $|G|^2$ growth of any Bursarial construction.
  \item \emph{Unbalanced weights cost some of the gain.} When the weights are
    unevenly spread over the residues, the busiest residue sets the length of
    $q$ (for example, $D_4$ gives $5.9\times$ rather than~8 for $S(4)$,
    $k = 2$).
  \item \emph{Exact further reductions.} The packing is built exactly, with
    coefficients in the cyclotomic integers $\mathbb{Z}[\zeta_L]$ for
    $L = \operatorname{lcm}(r, 4)$, which contain both $\zeta_r$ and~$i$. Each
    number is stored as a short integer vector (length $\phi(L)$: 2 for
    $L = 4$, 4 for $L = 12$, 8 for $L = 20$), and the floating-point coordinates
    used for evaluation are derived from these. With
    \texttt{reductions="safe"} two exact Level~1 searches run.
    \begin{itemize}
      \item \emph{The $c = 1$ block.} Each $E_{1,j}$ equals $(p_j S)\,w$, where
        $S = \sum_g M_g$ and $p_j$ is a weight-0 coordinate built from integer
        rows. Any that is a rational combination of the invariant outputs or of
        earlier ones is dropped.
      \item \emph{$2 \le c \le c_{\max}$.} The outputs are expanded exactly in
        $\mathbb{Z}[\zeta_L]$. As an exact check of the construction, every
        omitted coefficient is verified to vanish and every kept one to be
        real. Then each output $P$ contributes the rows
        $P, \zeta P, \ldots, \zeta^{\phi(L)-1} P$ as integer vectors (``restriction
        of scalars''). Rational linear relations among these rows are exactly
        the linear relations among the outputs with coefficients in
        $\mathbb{Q}(\zeta_L)$, so the integer relation search applies unchanged.
    \end{itemize}
    As predicted, little is found beyond the symmetry zeros above:
    $S(4)$, $k=2$ goes from 718 to 715 and $S(5)$, $k = 4$ from 13141 to 13139
    ($c_{\max} = 3$).
  \item \emph{No further gain from substitutions.} Scaling $x$, scaling $y$,
    reversing $x \mapsto 1/x$, and conjugation appear to be the only
    substitutions compatible with $\prod (y + q)$. So beyond about $2r$
    (perhaps $4r$), this route offers nothing more.
\end{itemize}

\section{Layer 3: removing redundancy}

There is no obviously best method for removing the redundancy in $\Phi$.
This section records what we are optimising and which reductions are safe.

\subsection{Design objectives}

\begin{enumerate}
  \item \emph{Injectivity on every input.} Target~(B) must hold for every $L$,
    including degenerate ones: coincident or parallel vectors, and $v_i$ that
    fail to span $\F^k$. Separation for generic $L$ alone is not acceptable.
    We expect sets that separate everywhere to be strictly larger than sets that
    separate only generic $L$. The extra outputs are needed only to resolve
    confusions on degenerate inputs. (For indefinite metrics, where some orbits
    are not closed, read ``every input'' as ``every orbit closure''. No
    continuous map can do more, and nothing more is needed; see
    Remark~\ref{rem:metric}.)

  \item \emph{Few outputs.} The number of outputs is to be minimised subject to
    (1). It can never be smaller than the dimension of the quotient, which for
    $n \ge k$ is $nk - \tfrac{1}{2}k(k-1)$. By Dym and Gortler, about
    twice the input dimension suffices if dense combinations are allowed, but
    see (3).

  \item \emph{Low complexity of each output, as a function of the raw inputs.}
    Separating invariants that have high degree and depend on delicate
    cancellations can be so ``spiky'' that, plotted against the components of
    the $v_i$, they look almost like white noise. Anything that later
    \emph{looks at} the embedding then cannot generalise from it. So the
    reduction in~(2) is balanced against keeping each output simple. Specifically:
    \begin{itemize}
      \item Complexity is measured by \emph{degree in the components of the
        $v_i$}, not by the size of the expression in $d$'s and $e$'s. Each $d$
        has degree $2$ and each $e$ has degree $k$. A finer measure is the
        multidegree, with one degree for each vector $v_i$. Every Gram,
        Plücker and mixed $d$--$e$ relation is homogeneous, so rewriting
        through them changes the expression but never the function, and never
        its degree. The growth in expression size from rewriting $e\,e\,e\,e$
        as $d$'s is therefore a problem of how outputs are displayed, not of
        how smooth they are.
      \item Degree does control smoothness. By a Bernstein-type inequality, a
        polynomial of degree $r$ restricted to the unit sphere has angular
        gradient at most $r$ times its maximum modulus there.
      \item Outputs should be \emph{homogeneous} in $L$, or better still
        multi-homogeneous, so that each one scales cleanly with the overall
        size of $L$. Splitting an output into its homogeneous parts never
        harms injectivity, because the parts determine the whole, although it
        does increase the count. The Bursarial $E_{c,j}$ are homogeneous in
        $L$ only when $k = 2$.
      \item A homogeneous output $f$ of degree $r$ can optionally be replaced by
        $\operatorname{sign}(f)\,|f|^{1/r}$. This is a bijection of
        $\mathbb{R}$, so injectivity is preserved, and every output then
        scales linearly with $L$. The cost is that the output is no longer
        differentiable where $f = 0$.
      \item Dense combinations that mix many monomials or degrees, such as
        generic linear projections, are disfavoured.
    \end{itemize}

  \item \emph{Evaluation cost} is a secondary concern and is expected to be
    modest.
\end{enumerate}

\subsection{Reductions that are always safe}

Removing an output that is given by an \emph{explicit, verifiable formula}
(here, a linear one) in terms of the outputs kept can never harm injectivity,
and the formula itself is the proof. The following reductions therefore need
no judgement:
\begin{itemize}
  \item if $P \subseteq G$, omit the $e$'s from $w$ from the outset;
  \item drop coefficients that vanish identically;
  \item drop exact duplicates and scalar multiples;
  \item drop any output in the linear span of those kept. This can be detected
    formally (in the $d$'s and $e$'s) or as functions of $L$, where the
    Gram, Plücker and mixed relations hold automatically.
\end{itemize}
If the reduced set still separates, then every discarded output is still
\emph{determined} by the outputs kept. Equal kept outputs imply equal $\MS$, so
a discarded output is a well-defined function of the kept ones. What changes
beyond the reductions above is that this function is generally not linear,
and not even polynomial. For example, $x$ is determined by $x^3$ on
$\mathbb{R}$, but only through a cube root. Discarded invariants typically
satisfy polynomial equations over the kept ones, and injectivity is what
selects the correct root. The safety of such a reduction therefore cannot be
read off from an identity among the outputs. It needs a certificate valid on
every input: either an exact algebraic computation (of the \emph{separating
variety}), or an explicit decoding algorithm, with case splits for degenerate
inputs, that is proved to recover the orbit of $L$ as in target~(B). Tests at
random points, such as Jacobian rank or counting preimages, can only propose
candidates or reject them.

\subsection{Levels of certification}\label{sec:levels}

``Certified'' means that there is evidence that an output set separates
orbits on \emph{every} input (target~(B)). That evidence comes in levels of
decreasing strength.
\begin{description}
  \item[Level 0: by construction.] Injectivity follows from a general theorem,
    and nothing needs to be checked for the particular $(n, k, G)$. Example:
    the unpruned Bursarial embedding (by unique factorisation).
  \item[Level 1: explicit identity.] A Level~0 set from which outputs have
    been removed, each removed output being an explicit formula in the kept
    ones (``$0$'', ``$=$ output~$7$'', ``$= \tfrac12\,\text{output}_3 -
    \text{output}_9$''). The code verifies each formula exactly, in rational
    arithmetic, and injectivity is inherited from Level~0. These are the
    reductions of Section~4.2.
  \item[Level 2: proof of separation.] The removed outputs are not formulas
    of the kept ones, but there is a proof that the kept set separates every
    input. The proof can take one of three forms:
    \begin{enumerate}
      \item[(a)] a theorem proved for a whole family of cases and recorded in
        these notes, so that the code need only check that the requested
        $(n, k, G)$ satisfies its hypotheses;
      \item[(b)] an exact computer-algebra computation for this particular
        $(n, k, G)$, such as showing that the separating variety is exactly
        the union of the graphs of the group action. The result is stored as
        an artefact that can be checked again;
      \item[(c)] a decoder that reconstructs the orbit of $L$ from the
        outputs, with a written proof, including case splits for degenerate
        inputs, that it succeeds on every input. The code implements the
        decoder, and tests exercise it.
    \end{enumerate}
    Forms (a) and (c) rest on proofs checked by people, and (b) rests on the
    correctness of the computer-algebra system. None is formally verified.
    ``Certified'' therefore means that a proof exists that we stand behind,
    not that the proof has been checked by machine.
  \item[Level 3: evidence only.] Tests at random points pass (full Jacobian
    rank, correct preimage counts, no confusable pairs found by search), but
    there is no proof. Generic injectivity is then plausible, but nothing
    reliable is known about degenerate inputs. Level~3 sets are
    \emph{uncertified}.
\end{description}

\subsection{Composing reductions}\label{sec:composition}

Certified reductions (Levels 0--2) compose. Each step leaves a set of outputs
that is still injective on $\MS$, so whatever a step discards remains
determined, in principle, by what it keeps. The order in which reductions are
applied therefore cannot affect correctness. It can only affect what later
reductions are able to find: an early, suboptimal reduction may make a better
one harder to discover, but it never breaks injectivity. Uncertified
(Level~3) candidates carry no such guarantee, and do not compose safely.

Because every step is recorded with its certificate, different orders can be
tried from scratch and compared, at least for small cases.

Some reductions prune outputs of a fixed embedding (for example, the safe
reductions). Others replace the layer~2 embedding itself by another that is
injective by construction (for example, involution packing, below). Later
reductions then act on the new outputs, and the chain stays certified.

\subsection{A pitfall: pairwise zips do not determine a table}\label{sec:pairwise}

A tempting way to exploit structure is to embed each row of the multiset
(viewed as a table whose columns are its members) separately, and then
``zip'' rows together in pairs, for example by embedding the multiset of
complex numbers $u + iv$ formed from two rows. \emph{This is not injective.} Two
different tables can agree on every pairwise zip. In the smallest example below,
the columns are the members:
\[
  T: \begin{array}{cccc} A & A & B & B \\ X & Y & X & Y \\ u & v & v & u \end{array}
  \qquad\qquad
  T': \begin{array}{cccc} A & A & B & B \\ X & Y & X & Y \\ v & u & u & v \end{array}
\]
Every pair of rows has the same multiset of column pairs in $T$ as in $T'$,
but $(A, X, u)$ is a column of $T$ and not of $T'$. A less symmetric example is
\[
  T: \begin{array}{ccccccc} A & A & B & B & B & C & C \\ X & Y & X & Y & Z & Y & Z \\
                              v & w & w & v & u & u & w \end{array}
  \qquad
  T': \begin{array}{ccccccc} A & A & B & B & B & C & C \\ X & Y & X & Y & Z & Y & Z \\
                               w & v & v & u & w & w & u \end{array}
\]
This counterexample stopped the development of the ``lumped multiset zip''
idea in the Set Encodings Redux notes.

$\MS$ is not a general table: it is a $G$-orbit. In an orbit, the values of a
slot orbit appear equally often, which neither example above respects. But on
degenerate inputs, coinciding values merge those counts into unequal ones. So
orbit structure cannot be assumed to rescue pairwise (or $r$-wise) zips. Any
scheme built on partial zips needs a proof, valid on every input, that uses the
orbit structure.

\section{Code interface (draft for discussion)}

This section specifies what code in Echinocoder should do, before any code is
written. It is intended to be iterated on.

\subsection{Two phases: build once, evaluate many times}

\paragraph{Build (once, may be slow).}
The user supplies a specification and receives an embedding function together
with its metadata:
\begin{verbatim}
    embed_fn, meta = build_invariant_embedding(k, group_spec, options)
\end{verbatim}
\begin{itemize}
  \item \texttt{k}: the dimension of each vector.
  \item \texttt{group\_spec}: defines $n$, the labels, and $G$. $G$ is a
    direct product of symmetric groups, one for each \emph{species}, for
    example $((a, b), (p))$ for $S(2) \times S(1)$, together with a flag for
    whether $P$ is included. This matches symatom's existing
    \texttt{VectorType}/\texttt{Context}, which already builds $w$ (as
    \texttt{repS}) and the signed-permutation action, and should be reused.
    Parity is not yet a group element in symatom. This class of groups is
    sufficient for now. Arbitrary permutation subgroups are out of scope, but
    the design should not rule them out.
  \item \texttt{options}: control the construction (see below).
\end{itemize}
Symbolic work happens only at build time: constructing $w$ and $G$, forming
the embedding, applying reductions, and recording provenance.

\paragraph{Evaluate (many times, must be cheap).}
The target framework is PyTorch, as used by the IWPC package
(\url{https://github.com/jjhwilkinson/iwpc}, which depends on \texttt{torch}
and \texttt{lightning}). \texttt{embed\_fn} is therefore a
\texttt{torch.nn.Module}:
\begin{verbatim}
    Y = embed_fn(X)     # X.shape == (..., n, k)  ->  Y.shape == (..., D)
\end{verbatim}
\begin{itemize}
  \item Any leading batch dimensions (``\texttt{...}'') are allowed, and may
    differ from call to call. $n$ and $k$ are fixed at build time, and inputs
    with any other trailing shape are rejected.
  \item Along the $n$ axis, \texttt{X[..., i, :]} is the vector with the
    $i$-th label of \texttt{group\_spec}. That ordering convention is part of
    the specification.
  \item Internally the pipeline is: \texttt{X} $\to$ the $d$'s and $e$'s
    ($w$, shape $(\ldots, m)$) $\to$ the polynomials chosen at build time
    $\to$ optional root transform $\to$ \texttt{Y}. Only tensor operations are
    used, so autograd works throughout. No group elements are enumerated at
    evaluation time. Any constants are registered as buffers, so that
    \texttt{.to(device)} and \texttt{state\_dict()} behave as usual.
  \item \texttt{Y} must separate orbits for \emph{every} input, and must be
    differentiable almost everywhere, so that it can be used for backprop.
    Differentiability may fail at corners: for example, the root transform
    is not differentiable where $f = 0$. Its gradient there must be defined
    explicitly (finite, rather than \texttt{inf} or \texttt{nan}).
  \item \emph{IWPC compatibility, without depending on IWPC.} Echinocoder must
    not have a hard dependency on IWPC. It should, however, fit naturally into
    IWPC, using the same packages (\texttt{torch} and \texttt{numpy}) and the
    same conventions as IWPC's \texttt{Encoding} base class:
    \begin{itemize}
      \item an \texttt{nn.Module} whose \texttt{input\_shape} and
        \texttt{output\_shape} are registered as integer buffers, and which acts
        on tensors of shape \texttt{(..., *input\_shape)};
      \item \texttt{input\_shape} $= (n, k)$ and \texttt{output\_shape} $= D$;
      \item a flattened variant with \texttt{input\_shape} $= nk$
        (vector-major: component $c$ of $v_i$ at position $ik + c$), because
        IWPC's \texttt{\&} concatenation needs vector inputs;
      \item NumPy-style docstrings, as in IWPC.
    \end{itemize}
    Wrapping the module as an IWPC \texttt{Encoding} should then be a matter
    of a few lines. Echinocoder may ship such an adaptor, provided IWPC is imported
    only when the adaptor is used, so that nothing else in Echinocoder needs
    IWPC installed.
  \item NumPy inputs may be accepted for convenience, for example in tests
    and documentation, by conversion to tensors.
\end{itemize}

\subsection{Options (initial list)}

\begin{itemize}
  \item \texttt{embedder}: which multiset embedding to use (layer 2), chosen
    from a registry so that others can be added. Currently
    \texttt{"bursar"} (Section~3.1), \texttt{"bursar\_involution"} (involution
    packing, Section~\ref{sec:involution}) and \texttt{"bursar\_graded"} (graded
    packing, Section~\ref{sec:graded}), all Level~0. Graded packing subsumes
    involution packing, which is kept as the simpler illustration.
  \item \texttt{reductions}: which reductions to apply, as an ordered list or
    a level. Early versions may support only \texttt{none}. Next would come
    the safe reductions of Section~4 (zero, duplicate, scalar multiple, linear
    span, which are Level~1; and dropping the $e$'s when $P \subseteq G$,
    which is Level~0 for the smaller $w$, because the $d$'s alone determine $L$
    up to $O(k)$).
    Level~2 reductions are applied only if the corresponding proof, artefact or
    decoder exists in the code. The main builder never returns a Level~3 set.
    See Section~\ref{sec:levels}.
  \item \texttt{homogenise}: split outputs into homogeneous (or
    multi-homogeneous) parts in $L$.
  \item \texttt{root\_normalise}: apply $\operatorname{sign}(f)|f|^{1/r}$ to
    each homogeneous output of degree $r$.
  \item \texttt{field}: $\mathbb{R}$ initially. Whether to support
    $\mathbb{C}$ is still open.
  \item \texttt{dtype}: the precision of evaluation, \texttt{float32} by
    default or \texttt{float64}. It affects only the evaluator. The build-time
    symbolic work is exact, with integer or rational coefficients. In keeping
    with the usual \texttt{torch} conventions, the module can also be
    converted afterwards with \texttt{.to(dtype)}.
\end{itemize}
An embedding that works in \texttt{float64} but not in \texttt{float32} is
already too fragile to be useful. Keeping outputs of low degree and low noise
(objective~3 of Section~4) is partly what makes \texttt{float32}
sufficient, so the designs should aim to work in \texttt{float32}. As a
diagnostic, the build could compare \texttt{float32} and \texttt{float64}
evaluations on random inputs and record the discrepancy for each output in
\texttt{meta}.

\subsection{Metadata}

\texttt{meta} is for documentation and auditing only. It is never used in
evaluation. It should be serialisable (for example as JSON), and should
contain or allow generation of:
\begin{itemize}
  \item the spec and options it was built from, plus the Echinocoder
    version or commit, so that the embedding can be rebuilt and checked;
  \item $D$, together with the slot order of $w$;
  \item for each output: its symbolic formula (see below), its degree and
    multidegree in $L$, whether it is homogeneous, any transform applied, and
    its provenance (for example, ``$E_{2,12}$ of the Bursarial embedding'');
  \item a log of the reductions: each discarded output, the reason, and the
    evidence. Examples are ``identically zero'', ``equal to output~$i$'',
    ``$= \sum_i \lambda_i\,(\text{output } i)$'', or a named certificate;
  \item the \emph{certification level} (Section~\ref{sec:levels}) and what it
    rests on: the theorem, the identities, the stored artefact, or the decoder
    together with its proof;
  \item a LaTeX table generator for the outputs.
\end{itemize}

\paragraph{Symbolic form of the outputs.}
Evaluation works from the components of $L$, but the symbolic description
never expands below the $d$'s and $e$'s. These are its leaves. Nobody reading
\texttt{meta} needs to see $E F - p_x q_x - p_y q_y - p_z q_z$ in place of a
dot product. Degrees in $L$ are still reported, counting $2$ for each $d$ and
$k$ for each $e$.

The description may also use \emph{nested definitions}, naming intermediate
quantities once and referring to them by name. For example,
\[
  \mathrm{enc}[23] = f\bigl(h(w)\bigr) + g\bigl(h(w)\bigr),
  \qquad \text{with } h(w) = \ldots
\]
Natural intermediates are $q_g(x)$, the elementary symmetric polynomials
$E_c$ and their coefficients $E_{c,j}$, and orbit sums. This applies to both
the machine-readable form (a set of named definitions, such as a DAG of
expressions) and the LaTeX output (a table of outputs plus a table of
definitions). Its purpose is to keep each displayed formula short. Whether to
use nesting, and how deeply, is a presentation choice. It does not affect
evaluation.

\subsection{Placement and reuse}

The code lives in a new Echinocoder subpackage next to symatom and symcoder.
It should not inherit from the existing \texttt{MultisetEmbedder}, which works
on one input at a time and has no group specification.

symatom already builds something close to $w$ (\texttt{repS}) and the
signed-permutation action of $G$ on it, and it should be reused where it
matches. It may do exactly what is needed here, or something subtly
different. Before any part of it is relied on, its behaviour must be checked
against the definitions of Section~1, in particular:
\begin{itemize}
  \item which atoms appear: all self-dots, cross-dots and $e$'s, and nothing
    else;
  \item the canonical index order, and the sign conventions for $e$'s under
    relabelling;
  \item parity. symatom's \texttt{Atom.sign} is known to conflate the sign
    from permuting arguments with the sign from spatial parity. This was
    harmless while there was no parity operation, but $P$ is now being added.
\end{itemize}
Tests should compare the $w$ and $g \cdot w$ obtained from symatom with an
independent, direct implementation of Section~1 on small cases. symcoder's
encoders are numerical, mostly based on sorting, and are not Bursarial. So
the main thing to reuse there is perhaps its design experience rather than
its code.

\subsection{Scale}

The initial goal is $k \le 5$ and $n \le 10$. Larger would be welcome. The
goal has strong consequences. For a single species with $n = 10$, $k = 5$:
\[
  m = \tbinom{10}{2} + 10 + \tbinom{10}{5} = 307, \qquad
  |G| = 10! \approx 3.6 \times 10^6 \;(\times 2 \text{ with } P), \qquad
  D_{\text{Bursar}} \approx 2 \times 10^{15},
\]
whereas the quotient dimension, which is the lower bound on the number of
outputs, is only $nk - \tfrac{1}{2}k(k-1) = 40$. Hence:
\begin{itemize}
  \item The unpruned Bursarial embedding is feasible only for small $|G|$.
    There it serves as a reference baseline, against which reductions can be
    validated.
  \item For large $G$, the build cannot enumerate $G$, and cannot first form
    the unpruned output and then prune it. It must construct the reduced
    outputs directly. One way is to work with orbits of monomials in $w$
    (orbit sums), described by the pattern of their labels rather than by
    listing group elements. symatom's orbit enumerators may help here.
  \item Evaluation should use tensor contractions, for example
    $\sum_{i,j} d_{ij}^2 = \operatorname{tr}(\mathbf{d}^2)$, rather than
    expanded lists of terms, whose length grows with orbit size.
\end{itemize}
Even $k \le 5$, $n \le 10$ may prove too big. It is a goal, not a
requirement.

\subsection{Uncertified (Level 3) candidates}

The main builder returns only sets at Levels~0--2. Level~3 candidates are
useful for research, for example to see how small output sets can get and to
suggest which sets are worth proving. They are available only through a
separate, explicitly named research entry point, \emph{not} through a flag
or option on the main builder, so that they cannot be picked up by accident.
Their \texttt{meta} states that they are uncertified (Level~3), and lists the
evidence collected.

\section{Exploration log}\label{sec:exploration}

This section records what has been observed by running the code on very simple
cases. It also lists lines of enquiry those observations suggest. The lines of
enquiry are \emph{not} a plan. They are unordered, none is singled out as the
next thing to investigate, and other ideas are expected to be added and
compared with them.

\subsection{Tools and outputs}

\begin{itemize}
  \item \path{symembed/explore/simple_cases.py} writes each case out in full
    (\path{docs/explorations/simple_cases.pdf}, with every expansion in
    \path{simple_cases_full.txt}). It shows the slots of $w$, the action of
    $G$, every $E_{c,j}$ expanded in the $d$'s and $e$'s, and the effect of the
    safe reductions.
  \item The polynomial-closure tool, \path{symembed/explore/polynomial_closure.py}
    (report: \path{docs/explorations/polynomial_closure.txt}), visits outputs
    in order of $c$, then $j$. It keeps an output unless it is an exact
    polynomial in outputs already kept, and then retries each kept output
    against the others (``backward pruning''). Membership is tested either
    formally in $w$ or as functions of $L$, where the Gram and Plücker relations
    hold automatically. Identities are found modulo four primes, reconstructed
    by the Chinese remainder theorem, and verified exactly. Every kept set is
    therefore a Level~1 certified separating set: each Bursarial output is an
    explicit polynomial in it. Each case is run for both the plain and the
    involution-packed embedding.
\end{itemize}

\subsection{Observations}

The table gives the number of outputs in each case (all without parity).

\begin{center}
\begin{tabular}{lrrrrr}
  case & Bursarial & linear & poly.\ in $w$ & poly.\ in $L$ & quot.\ dim. \\ \hline
  $S(2)$, $(a,b)$, $k=1$             & 14  &     & 6  & 2  & 2 \\
  $S(3)$, $(a,b,c)$, $k=1$           & 174 &     & 21 & 5  & 3 \\
  $S(2)\times S(1)$, $(a,b,p)$, $k=1$ & 26  &     & 12 & 3  & 3 \\
  $S(2)$, $(a,b)$, $k=2$             & 11  & 8   & 5  & 4  & 3 \\
  $S(3)$, $(a,b,c)$, $k=2$           & 174 & 155 & 32 & 13 & 5 \\
  $S(2)\times S(1)$, $(a,b,p)$, $k=2$ & 26  & 22  & 15 & 11 & 5
\end{tabular}
\end{center}

\begin{enumerate}
  \item \emph{Polynomial identities do much more than linear ones}, and allowing
    the Gram and Plücker relations (working in $L$ rather than formally in $w$)
    does much more again.
  \item \emph{Low $c$ suffices for generation.} For $S(3)$ every kept output has
    $c \le 4$ ($k = 2$) or $c \le 3$ ($k = 1$), although the Bursarial
    embedding goes up to $c = |G| = 6$. All higher-$c$ outputs are polynomials
    in lower ones.
  \item \emph{For $k = 1$} (the classical symmetric-polynomial setting), the
    count reaches the quotient dimension for $S(2)$ and $S(2) \times S(1)$.
    For $S(3)$ it stops at~5. The search can only keep actual Bursarial outputs,
    and it is greedy, so this says nothing about how small a set could be.
  \item \emph{Polynomial closure gives generating sets, not minimal separating
    sets.} For $k = 2$ the counts stay well above the quotient dimension (11
    against 5, and 13 against 5).
  \item \emph{$S(2)$, $k = 2$, as a zip.} Here
    \[
      \MS = \{(d_{aa}, d_{bb}, d_{ab}, e_{ab}),\ (d_{bb}, d_{aa}, d_{ab}, -e_{ab})\}.
    \]
    Store the invariant $d_{ab}$ as it is, and zip $u + iv$ with
    $(u, v) = (d_{aa}, e_{ab})$ in one member and $(d_{bb}, -e_{ab})$ in the
    other. This gives
    $d_{aa} + d_{bb}$ (the imaginary part vanishes) and
    $(d_{aa} d_{bb} + e_{ab}^2) + i\, e_{ab}(d_{bb} - d_{aa})$: 4 reals in all.
    These coincide with the 4 kept by the polynomial closure, up to a
    polynomial change of variables. One of the identities needs the Gram
    relation $e_{ab}^2 = d_{aa}d_{bb} - d_{ab}^2$. The quotient dimension is
    3. For fixed $d_{aa} + d_{bb}$ and $d_{ab}$, the remaining freedom is a
    closed loop, so any 3-output scheme would have to mix coordinates
    nonlinearly.
  \item \emph{Complex structure for $k = 2$.} Identify each $v_i$ with a complex
    number $z_i$. Then
    \[
      \bar z_i z_j = d_{ij} + i\, e_{ij},
    \]
    so the $d$'s and $e$'s are the real and imaginary parts of products that
    are invariant under rotation.
  \item \emph{Reference covariants reach the quotient dimension generically.}
    \begin{itemize}
      \item For $S(2) \times S(1)$, $(a, b, p)$, $k = 2$: put
        $\alpha = \bar p a = d_{ap} + i e_{ap}$ and
        $\beta = \bar p b = d_{bp} + i e_{bp}$. When $p \ne 0$ the orbit is
        determined by $|p|^2 = d_{pp}$ and the unordered pair
        $\{\alpha, \beta\}$, and so by $d_{pp}$, $\alpha + \beta$ and
        $\alpha\beta$. That is 5 reals, equal to the quotient dimension, each
        of degree~2 in $w$.
        \begin{itemize}
          \item When $p = 0$, $\alpha = \beta = 0$ carries nothing. Adding the 4
            outputs of $S(2)$, $k = 2$ for $(a, b)$ covers that stratum, and
            $d_{pp}$ shows which case applies. That gives 9 outputs that
            separate every input, by a short decoder argument (Level~2c),
            against 11 from the polynomial closure.
        \end{itemize}
      \item For $S(3)$, $(a, b, c)$, $k = 2$: the sum $s = \sum_i z_i$ gives
        $\alpha_j = \bar s z_j = \sum_i (d_{ij} + i e_{ij})$. The elementary
        symmetric polynomials of the $\alpha_j$ give $1 + 2 + 2 = 5$ reals,
        equal to the quotient dimension, but only where $s \ne 0$. The stratum
        $s = 0$ needs something else.
      \item In both examples, each degenerate stratum on which the reference
        vanishes costs extra outputs. This is consistent with the expectation
        that sets separating every input are larger than sets separating only
        generic inputs.
    \end{itemize}
\end{enumerate}

\paragraph{Involution packing under reduction.}
The counts below run from outputs before reduction, to those kept by the
polynomial closure in $w$, to those kept by the polynomial closure in $L$.

\begin{center}
\begin{tabular}{lll}
  case & plain & involution-packed \\ \hline
  $S(2)$, $k=2$             & $11 \to 5 \to 4$    & $5 \to 5 \to 4$ \\
  $S(3)$, $k=2$             & $174 \to 32 \to 13$ & $90 \to 58 \to 13$ \\
  $S(2)\times S(1)$, $k=2$  & $26 \to 15 \to 11$  & $14 \to 12 \to 11$ \\
  $S(2)$, $k=1$             & $14 \to 6 \to 2$    & $8 \to 6 \to 2$ \\
  $S(3)$, $k=1$             & $174 \to 21 \to 5$  & $111 \to 24 \to 5$ \\
  $S(2)\times S(1)$, $k=1$  & $26 \to 12 \to 3$   & $17 \to 11 \to 3$
\end{tabular}
\end{center}

After full closure in $L$ the two agree in every case. Formally in $w$ the
packed embedding is sometimes worse (58 against 32 for $S(3)$, $k=2$). Packing
halves the raw count at no cost, but leaves unchanged what survives the
polynomial closure, which is itself only feasible for small cases.

\paragraph{Parity of graded outputs.}
Each output can be classified by what happens under $e \mapsto -e$, the action
of parity: even if unchanged, odd if it changes sign, mixed if neither. Testing
numerically, as polynomials in $w$, with an exact cross-check for low $c$:
\begin{center}
\begin{tabular}{lrrrr}
  case & outputs & even & odd & mixed \\ \hline
  $S(3)\times S(2)$, $k = 3$ & 191 & 17 & 6 & 168 \\
  $S(3)\times S(2)$, $k = 4$ & 124 & 19 & 1 & 104
\end{tabular}
\end{center}
Most outputs are mixed, because the packing combines $d$'s and $e$'s. If
outputs of definite parity are wanted, the simplest route is to split each
mixed output $o$ afterwards into $\tfrac12(o(L) + o(pL))$ and
$\tfrac12(o(L) - o(pL))$, at a cost of one extra output per mixed one. That
gives $191 + 168 = 359$ and $124 + 104 = 228$. A ``parity-graded'' packing,
using coordinates of pure parity placed so that the parity of each position
matches the parity of its coordinate, makes every coefficient pure
automatically: $E_{c,j}(p\cdot w) = (-1)^j E_{c,j}(w)$. But it does
\emph{worse} here, 395 and 273. The parity condition rules out the best $h$,
whose order $6$ is even and conflicts with it, and forces extra padding.

\paragraph{Comparison with hand-crafted separating sets.}
For five Lorentz vectors with $G = S(3) \times S(2)$ ($k = 4$),
\cite{GHL20} (Table~3) gives 26 parity-even variables and \cite{LHG22}
(conclusions) 28 parity-odd ones. Together they form a separating set of 54
variables. The parity-odd part was crafted by arguments that could not be
automated. Those sets are restricted to timelike or null vectors, whereas the
constructions here are injective on $\MS$ for all inputs, so the comparison
flatters the crafted sets somewhat. The fully automatic graded construction
(with the safe reductions) gives:
\begin{center}
\begin{tabular}{lrrr}
  & crafted & here & ratio \\ \hline
  with $P$ in $G$ (parity-even only)          & 26 & 112 & 4.3 \\
  without $P$, outputs as they come             & 54 & 124 & 2.3 \\
  without $P$, split into definite parity       & 54 & 228 & 4.2
\end{tabular}
\end{center}
The quotient dimension is 14. Nonlinear (polynomial) identities and the Gram
and Pl\"ucker relations have not yet been exploited on graded outputs. On the
small cases those were the largest reductions of all.

\subsection{Possible lines of enquiry (unordered, not exhaustive)}

\begin{itemize}
  \item The order of quotienting. The orbit $\MS$ has $|G| = n!$ members only
    because rotations are removed first (via the $d$'s and $e$'s) and the
    permutations afterwards. Removing the permutations first, by embedding the
    multiset of the $n$ raw vectors, costs only about $k n^2/2$ outputs. But
    the result is then covariant, not invariant, under rotations, and the
    difficulty moves into rotation invariants of it. For $k = 2$ these are
    phases: the $c$-th elementary symmetric function of the $z_i$ has weight
    $c$. That route would need Level~2 arguments for degenerate inputs. The
    point-cloud literature has completeness theorems (for example
    simplex-wise distance distributions, and $(d-1)$-WL tests for
    $d$-dimensional clouds) that might supply them.
  \item Extending the safe reductions from linear to polynomial identities
    (Level~1). On the cases above this alone takes, for example, 26 down to
    11.
  \item Reference covariants, for $k = 2$ via the complex structure, with the
    degenerate strata handled by further covariants or by recursion to smaller
    cases. Open questions: whether there is an analogue for $k \ge 3$, which
    has no complex structure, and how the number of strata grows.
  \item Understanding why low $c$ suffices for generation, and whether there is
    a useful bound on the $c$ needed in terms of the structure of $G$.
  \item Further simple cases for comparison: parity, $k = 3$, and two
    non-trivial species.
  \item Complex packing of $w$: pair the real slots of $w$ into
    $\lceil m/2 \rceil$ complex numbers, embed the multiset of packed vectors
    with the complex Bursarial embedding, and write out real and imaginary
    parts. The raw count changes very little, because the Bursarial size is
    linear in $m$. For even $m$ it saves $|G|(|G|-1)/2$ reals (as already
    noted in the legacy embedder's docstring), which is a fraction of
    about $1/m$. For odd $m$, padding with zero costs $|G|$ reals. Any real
    benefit would have to show up in what survives reduction, and that may
    depend on the pairing (the order of the slots). A factor of~2 would need an
    element of $G$ acting on the packed vector as complex conjugation, and
    without parity no permutation other than the identity does so.
  \item Involution packing, which generalises the $z \mapsto i\bar z$ trick and
    does not need parity. Choose any involution $g_0 \in G$. It splits $w$
    into a $+1$ part and a $-1$ part. A slot pair $(s, t)$ swapped by $g_0$
    contributes $w_s + w_t$ and $w_s - w_t$, with the signs adjusted as
    required, and a fixed slot contributes to the part given by its sign. Pack
    the $+1$ coordinates as real parts and the $-1$ coordinates as imaginary
    parts, padding with zeros if the counts differ. Then $g_0$ acts on every
    packed vector as complex conjugation. The multiset of packed vectors is
    therefore closed under conjugation, so $\prod_g (y + q_g(x))$ has real
    coefficients, and only those need storing.
    \begin{itemize}
      \item This is Level~0: the complex Bursarial embedding is injective by
        unique factorisation, and the packing is a bijection.
      \item The number of packed slots is
        $m'' = (m + |\operatorname{tr} g_0|)/2$, so the involution with
        smallest $|\operatorname{tr}|$ is best.
      \item Raw counts fall to roughly half: $11 \to 5$ for $S(2)$, $k=2$;
        $174 \to 90$ for $S(3)$, $k=2$; $26 \to 14$ for $S(2)\times S(1)$,
        $k=2$; $138060 \to 79980$ for $S(5)$, $k=4$; $1884 \to 948$ for
        $S(3)\times S(2)$, $k=3$.
      \item Now implemented (Section~\ref{sec:involution}, embedder
        \texttt{bursar\_involution}). What survives full polynomial closure is
        unchanged on the simple cases (see above). The factor of~2 has since
        been exceeded by graded packing (Section~\ref{sec:graded}).
    \end{itemize}
  \item Graded packing (Section~\ref{sec:graded}, embedder
    \texttt{bursar\_graded}): a dihedral subgroup acting on $x$ by roots of
    unity and on the coefficients by conjugation. It gives $5$--$10\times$ on
    the cases tried ($S(5)$, $k = 4$: $138060 \to 13141$). Open questions:
    \begin{itemize}
      \item what survives polynomial closure;
      \item exact Level~1 reductions over $\mathbb{Q}(\zeta_r)$;
      \item whether other structures of $G$ can be exploited in the same way.
      \item Going beyond a factor of~2 with two commuting involutions would
        need an algebra with two independent conjugations. Bicomplex numbers
        have zero divisors, which breaks unique factorisation, and quaternions
        do not commute, so the product depends on the order of its factors.
    \end{itemize}
\end{itemize}

\begin{thebibliography}{9}
\bibitem{GHL20} B.~Gripaios, W.~Haddadin and C.~G.~Lester, \emph{Lorentz and
  permutation invariants of particles I}, arXiv:2003.05487,
  doi:10.1088/1751-8121/abe58c.
\bibitem{LHG22} C.~G.~Lester, W.~Haddadin and B.~Gripaios, \emph{Lorentz and
  permutation invariants of particles III: constraining non-standard sources of
  parity violation}, Int.\ J.\ Mod.\ Phys.\ A \textbf{37} (2022) 2250093,
  arXiv:2008.05206.
\end{thebibliography}

\end{document}
